\documentclass[12pt]{article}
\usepackage{amsmath}
\usepackage{amssymb}
\usepackage{float}
\usepackage{graphicx}
\usepackage[left=2cm, right=2cm]{geometry}
\usepackage{titlesec}
\usepackage{svg}

\title{\vspace{-3cm}Simulation of the Noisy Harmonic Oscillator:\\
Gaussian and L\'{e}vy White Noise Forcing}
\author{Tristan Mihocko}
\date{May 2026}

\titlespacing{\subsection}{2em}{3.25ex}{1.5ex}

\begin{document}

\maketitle

\section*{Abstract}
The harmonic oscillator is one of the most fundamental models in applied
mathematics, describing physical systems from mechanical springs and electrical
circuits to molecular vibrations. In idealized settings the system evolves
deterministically, but in practice external forces are often random: thermal
fluctuations, turbulent flow, or unpredictable impulses perturb the trajectory
continuously. Replacing the deterministic forcing with white noise converts
the governing ordinary differential equation into a stochastic differential
equation (SDE), whose solutions are random processes rather than fixed curves.

This project presents an interactive simulation of the noisy harmonic oscillator
driven by either Gaussian white noise or L\'{e}vy stable white noise. The
simulation numerically integrates the SDE using the Euler--Maruyama method and
produces two plots: the displacement time series $x(t)$ and the phase portrait
in the $(x, v)$ plane. Both plots update in real time through interactive
\texttt{matplotlib} sliders, allowing the user to explore how the physical
parameters and the character of the noise shape the system's behavior. The
tool also compares the qualitatively different trajectories produced by Gaussian
and L\'{e}vy noise, illustrating the role of the stability index $\alpha$ in
generating heavy-tailed, jump-dominated dynamics.

\section{Mathematical Background}

This section develops the mathematical framework underlying the simulation,
beginning with the classical oscillator and building up to the stochastic
setting and the two noise models.

\subsection{The Classical Damped Harmonic Oscillator}

The damped harmonic oscillator is governed by the second-order linear ODE:

\begin{equation}
    m\,x''(t) + c\,x'(t) + k\,x(t) = 0
    \label{eq:ode}
\end{equation}

where $x(t)$ is the displacement from equilibrium, $m > 0$ is the mass,
$c\ge0$ is the damping coefficient, and $k>0$ is the spring constant.
Two derived values characterize the system's behavior: the natural frequency at which the undamped system oscillates \eqref{eq:frequency}, and the critical damping value \eqref{eq:damping}, which separates oscillatory from non-oscillatory behavior. 

\begin{equation}
    \omega_0 = \sqrt{k/m} 
    \label{eq:frequency}
\end{equation}

\begin{equation}
    c_{\text{crit}} = 2\sqrt{km}
    \label{eq:damping}
\end{equation}

    
The qualitative character of the solution depends on the ratio
$\zeta = c / c_{\text{crit}}$:

\begin{itemize}
    \item \textbf{Underdamped} ($c < c_{\text{crit}}$): the system oscillates
          at frequency $\omega_d = \omega_0\sqrt{1-\zeta^2}$ with amplitude
          decaying as $e^{-\zeta\omega_0 t}$. In the phase portrait this
          appears as an inward spiral converging to the origin.
    \item \textbf{Critically damped} ($c = c_{\text{crit}}$): the system
          returns to equilibrium as quickly as possible without oscillating.
          The phase portrait shows a trajectory curving directly into the origin.
    \item \textbf{Overdamped} ($c > c_{\text{crit}}$): the system returns to
          equilibrium slowly without oscillating, as a sum of two decaying
          exponentials. The phase portrait shows a trajectory that approaches
          the origin along a nearly straight path.
\end{itemize}

\subsection{The Phase Portrait}

The phase portrait is a plot of the velocity $v(t) = x'(t)$ and the displacement $x(t)$, tracing the trajectory of the system in the two-dimensional space $(x, v)$. Rather than showing how $x$ or $v$ evolve over time individually, the phase portrait shows the relationship between 
position and velocity across the entire trajectory, making it easier for the human eye to notice qualitative features.

For the unforced damped harmonic oscillator, the phase portrait is a smooth spiral
(underdamped) or a curve approaching the origin (overdamped or critically
damped). Adding noise disrupts this smooth structure: the trajectory wanders
randomly through the $(x, v)$ plane, tracing an irregular, space-filling
curve rather than a clean spiral. Under L\'{e}vy noise, sudden large
increments in $v$ produce sharp, nearly vertical, jumps in the phase portrait,
which are very distinct from the gradual drift seen under Gaussian noise.

\subsection{Stochastic Processes}

A stochastic process is a collection of random variables $\{X_t\}_{t\ge0}$ indexed by time, representing a probability distribution over infinitely many possible trajectories called \textit{sample paths}. A single run of the simulation produces one such path. Re-running the simulation with the same parameters produces a different path with the same statistical properties. The randomness is not a numerical error; it is a core feature of the model.

\subsection{Stochastic Differential Equations}

While an ODE describes a system whose future state is completely determined by
its current state via deterministic rules, an SDE introduces randomness
directly into those rules by adding a noise term which is a stochastic process. 
Consequently, the solution of an SDE is not a single curve but a stochastic
process.

To add stochastic forcing to the classic damped harmonic oscillator \eqref{eq:ode}, we replace the right-hand side with a white noise term $F(t)$:

\begin{equation}
    m\,x''(t) + c\,x'(t) + k\,x(t) = F(t)
    \label{eq:sde}
\end{equation}

Introducing the velocity $v(t) = x'(t)$ reduces this to a first-order system,
which is the standard form for numerical integration:

\begin{align}
    dx &= v\,dt \label{eq:sys1} \\
    dv &= \left(-\frac{c}{m}\,v - \frac{k}{m}\,x\right)dt
          + \frac{1}{m}\,dW_t \label{eq:sys2}
\end{align}

where $dW_t$ is the differential noise increment over the time step $dt$.
Its distribution depends on the chosen noise model.

\subsection{Gaussian White Noise}

The standard model of continuous random forcing is the Wiener process
(Brownian motion) $W_t$. Its increments $dW_t = W_{t+dt} - W_t$ are:

\begin{itemize}
    \item independent for non-overlapping time intervals,
    \item normally distributed: $dW_t \sim \mathcal{N}(0,\, \sigma^2 dt)$,
    \item continuous almost surely.
\end{itemize}

The parameter $\sigma > 0$ is the noise scale. Because the variance of each
increment is proportional to $dt$, the typical magnitude of a fluctuation
over a small step scales like $\sigma\sqrt{dt}$, which is much larger than
the drift contribution, which scales linearly with $dt$, for small $\Delta t$. This is what causes the famous roughness of Brownian paths.
The resulting SDE is treated in the framework of It\^{o} calculus, with It\^{o}'s lemma providing the fundamental rules for differentiation and analysis of stochastic processes.

\subsection{L\'{e}vy Stable White Noise}

A more general family of noise processes is built from the L\'{e}vy stable
distributions. A random variable $X$ follows the stable distribution
$S(\alpha, \beta, \sigma, \mu)$ characterized by four parameters:

\begin{itemize}
    \item $\alpha \in (0, 2]$: the \textbf{stability index}. Controls how
        heavy the tails of the distribution are. When $\alpha = 2$ the
        distribution is exactly Gaussian. When $\alpha < 2$ the tails of the distribution satisfy $f(x) \sim |x|^{-(\alpha+1)}$, meaning the distribution
        has infinite variance and the noise can produce arbitrarily large
        increments with a non-negligible probability.
    \item $\beta \in [-1, 1]$: the \textbf{skewness parameter}. When $\beta = 0$
        the distribution is symmetric. When $\beta = 1$ it is completely
        right-skewed and when $\beta = -1$ completely left-skewed.
    \item $\sigma > 0$: the \textbf{scale parameter}, controlling the spread
        of the distribution.
    \item $\mu \in \mathbb{R}$: the \textbf{location parameter} shifts the distribution along the real line. It is standard practice to set $\mu = 0$, centering noise about the origin.
\end{itemize}

L\'{e}vy stable distributions are defined through their characteristic function
rather than a closed-form density.

\begin{equation}
    \log \mathbb{E}\!\left[e^{itX}\right] =
    -\sigma^\alpha |t|^\alpha
    \!\left(1 - i\beta\,\mathrm{sgn}(t)\Phi(\alpha, t)\right)
    + i\mu t
    \label{eq:levy_stable}
\end{equation}

\begin{equation}
    \Phi(\alpha, t) = \begin{cases}
        \tan \left(\frac{\pi\alpha}{2}\right) & \text{if } \alpha \ne 1 \\
        -\frac{2}{\pi}\ln|t| & \text{if } \alpha = 1
    \end{cases}
\end{equation}

The L\'{e}vy stable noise increment over a step $\Delta t$ is drawn from:

\begin{equation}
    \Delta W \sim S\!\left(\alpha,\, \beta,\,
    \sigma\,\Delta t^{1/\alpha},\, 0\right)
\end{equation}

The $\Delta t^{1/\alpha}$ scaling is the natural generalization of the
$\sqrt{\Delta t} = \Delta t^{1/2}$ scaling of Gaussian increments, recovering
the Gaussian case when $\alpha = 2$.
For $\alpha < 2$ the distribution has heavier tails than the Gaussian, producing
occasional large jumps in $x(t)$ and sharp excursions in the phase portrait.
This makes L\'{e}vy noise a natural model for systems subject to rare but
extreme shocks.

\section{Numerical Method}

\subsection{The Euler--Maruyama Scheme}

The Euler--Maruyama method is the SDE analog of the famous Euler
method for ODEs. Applied to the system \eqref{eq:sys1}--\eqref{eq:sys2}
with time step $\Delta t$, it gives the update:

\begin{align}
    x_{i+1} &= x_i + v_i\,\Delta t \label{eq:em1} \\
    v_{i+1} &= v_i
              + \left(-\frac{c}{m}\,v_i - \frac{k}{m}\,x_i\right)\!\Delta t
              + \frac{\Delta W_i}{m} \label{eq:em2}
\end{align}

where $\Delta W_i$ is a single random draw from the noise distribution.
For Gaussian noise, $\Delta W_i \sim \mathcal{N}(0, \sigma^2\Delta t)$.
For L\'{e}vy noise,
$\Delta W_i \sim S(\alpha, \beta, \sigma\,\Delta t^{1/\alpha}, 0)$.

Equation \eqref{eq:em1} advances the position using the current velocity,
and equation \eqref{eq:em2} advances the velocity using the current
deterministic drift plus a random noise kick. The key difference from the
ODE Euler method is the $\Delta W_i$ term, which injects randomness
at every step.

\subsection{Convergence}

For SDEs driven by Wiener processes, the Euler--Maruyama method has strong
convergence of order $1/2$: the expected absolute error
$\mathbb{E}[|x(T) - x_N|]$ scales as $O(\sqrt{\Delta t})$.
This is slower than the $O(\Delta t)$ convergence of forward Euler for ODEs,
a consequence of the nowhere-differentiability of Brownian paths.
For L\'{e}vy-driven SDEs with $\alpha < 2$ the convergence theory is more
involved due to the infinite variance of the increments; in practice the
step size $\Delta t = 0.01$ produces visually stable trajectories for the
parameter ranges used here.

\subsection{Numerical Parameters}

\begin{table}[H]
    \centering
    \begin{tabular}{lll}
        \hline
        \textbf{Parameter}    & \textbf{Symbol}  & \textbf{Value} \\
        \hline
        Mass                  & $m$              & $1.0$ kg \\
        Time step             & $\Delta t$       & $0.01$ s \\
        Total time            & $T$              & $40.0$ s \\
        Number of steps       & $N$              & $4000$ \\
        Initial displacement  & $x(0)$           & $0.5$ m \\
        Initial velocity      & $v(0)$           & $0.0$ m/s \\
        \hline
    \end{tabular}
    \caption{Fixed numerical parameters used in all simulations.}
    \label{tab:num_params}
\end{table}

\section{Implementation}

The simulation is implemented in Python using \texttt{numpy},
\texttt{scipy}, and \texttt{matplotlib}.

\subsection{Noise Sampling}

At each time step a noise increment is drawn from the selected distribution.
For Gaussian noise, \texttt{numpy.random.normal(0, sigma * sqrt(dt))} is used.
For L\'{e}vy stable noise, \texttt{scipy.stats.levy\_stable.rvs} is called
with the user-specified $\alpha$ and $\beta$ and scale $\sigma\,\Delta t^{1/\alpha}$.

\subsection{Visualization}

The figure contains two panels that update together whenever any parameter is changed.

\textbf{Left panel --- displacement $x(t)$:} plots the displacement against time
over the full $40$ s simulation window. This shows the oscillatory structure of
the trajectory, the effect of damping on amplitude, and the character of the noise
(smooth fluctuations for Gaussian, intermittent jumps for L\'{e}vy).

\textbf{Right panel --- phase portrait:} plots the velocity $v(t)$ against the
displacement $x(t)$ as a parametric curve, tracing the system's trajectory. For the unforced underdamped oscillator this is a smooth inward spiral. Adding noise disrupts this into an irregular, wandering curve. The white
marker indicates the starting point $(x(0), v(0)) = (0.5, 0)$.

\subsection{Interactive Controls}

Five \texttt{Slider} widgets allow real-time adjustment of the simulation
parameters:

\begin{table}[H]
    \centering
    \begin{tabular}{llll}
        \hline
        \textbf{Slider} & \textbf{Parameter} & \textbf{Range} & \textbf{Default} \\
        \hline
        Damping $c$           & Damping coefficient & $[0.01,\;3.0]$  & $0.20$ \\
        Spring $k$            & Spring constant     & $[0.50,\;12.0]$ & $4.00$ \\
        Noise scale $\sigma$  & Noise intensity     & $[0.01,\;2.0]$  & $0.30$ \\
        L\'{e}vy $\alpha$     & Stability index     & $[0.50,\;2.0]$  & $1.50$ \\
        L\'{e}vy $\beta$      & Skewness            & $[-1.0,\;1.0]$  & $0.00$ \\
        \hline
    \end{tabular}
    \caption{Interactive sliders, their parameters, ranges, and default values.}
    \label{tab:sliders}
\end{table}

A \texttt{RadioButtons} widget switches between L\'{e}vy stable and Gaussian
noise.
A \textbf{Re-run} button generates a new noise path at the current
parameter values.
A text panel below the controls displays the natural frequency $f_0$,
critical damping $c_{\text{crit}}$, and current damping regime, updating
automatically as the sliders are moved.

\section{Results}

\subsection{Unforced System ($\sigma \approx 0$)}

Setting the noise scale near zero yields approximately deterministic behavior.
With $c = 0.2$ and $k = 4.0$ (underdamped, $c_{\text{crit}} = 4.0$), the
$x(t)$ plot shows an oscillation at $f_0 \approx 0.318$ Hz whose
amplitude decays exponentially.
The phase portrait traces a clean inward spiral converging to the origin.
Increasing $c$ past $c_{\text{crit}}$ eliminates the oscillation: $x(t)$
decreases monotonically and the phase portrait shows a curve approaching
the origin along an almost straight path.

\begin{figure}[H]
    \centering
        \begin{minipage}{0.48\textwidth}
        \centering
        \includesvg[width=\textwidth]{levy_displacement_unforced.svg}
        \caption{Displacement of oscillator under L\'evy noise with noise scale $\sigma\approx 0$}
    \end{minipage}
    \hfill
    \begin{minipage}{0.48\textwidth}
        \centering
        \includesvg[width=\textwidth]{levy_phase_unforced.svg}
        \caption{Phase portrait of oscillator under L\'evy noise with noise scale $\sigma\approx 0$}
    \end{minipage}
\end{figure}

\subsection{Gaussian Noise}

Under Gaussian forcing with $\sigma = 0.3$, the $x(t)$ plot shows sustained
oscillations near $\omega_0$ coupled with small random perturbations.
The amplitude no longer decays to zero. Instead, the noise continuously
re-excites the oscillator. In the phase portrait the clean spiral is replaced
by a fuzzy, irregular band tracing approximately the same elliptical region,
centered on the origin. Increasing $\sigma$ widens this band; increasing $c$
shrinks it by suppressing the oscillations.
Different noise realizations at the same parameter values look broadly similar,
reflecting the finite variance of the Gaussian distribution.

\begin{figure}[H]
    \centering
        \begin{minipage}{0.48\textwidth}
        \centering
        \includesvg[width=\textwidth]{gaussian_displacement.svg}
        \caption{Displacement of oscillator under Gaussian noise}
    \end{minipage}
    \hfill
    \begin{minipage}{0.48\textwidth}
        \centering
        \includesvg[width=\textwidth]{gaussian_phase.svg}
        \caption{Phase portrait of oscillator under Gaussian noise}
    \end{minipage}
\end{figure}


\subsection{L\'{e}vy Noise}

Switching to L\'{e}vy stable noise with $\alpha = 1.5$, $\beta = 0$, and
$\sigma = 0.3$ produces qualitatively different behavior in both plots.
The $x(t)$ time series is punctuated by intermittent large-amplitude jumps
entirely absent under Gaussian noise, where the displacement suddenly shifts
to a value far from its current level before the oscillator gradually rings
back. In the phase portrait, these events appear as sudden near-vertical
excursions in the $v$ direction, representing instantaneous large velocity
kicks, followed by curved segments as the system oscillates back.

\begin{figure}[H]
    \centering
        \begin{minipage}{0.48\textwidth}
        \centering
        \includesvg[width=\textwidth]{levy_displacement.svg}
        \caption{Displacement of oscillator under L\'evy noise}
    \end{minipage}
    \hfill
    \begin{minipage}{0.48\textwidth}
        \centering
        \includesvg[width=\textwidth]{levy_phase.svg}
        \caption{Phase portrait of oscillator under L\'evy noise}
    \end{minipage}
\end{figure}

Reducing $\alpha$ toward $0.5$ increases both the frequency and magnitude of
jumps. At $\alpha = 1.0$ the trajectory is dominated by jump events and the
underlying oscillatory structure is barely visible between them. The skewness
$\beta$ controls the direction of the jumps: at $\beta = 1$ the jumps
are predominantly in the positive $x$ direction, which is visible in the
phase portrait as asymmetric excursions above and to the right of the origin.

\subsection{Comparison of Noise Models}

The most direct way to compare the models is to switch between them using
the radio buttons at identical parameter values. The key qualitative
differences are summarized in the table below.

\begin{table}[H]
    \centering
    \begin{tabular}{lll}
        \hline
        \textbf{Feature} & \textbf{Gaussian ($\alpha=2$)} & \textbf{L\'{e}vy ($\alpha < 2$)} \\
        \hline
        Increment distribution & Finite variance & Infinite variance \\
        Trajectory character   & Smooth, small fluctuations & Intermittent large jumps \\
        Phase portrait         & Fuzzy band / irregular spiral & Spiral with sharp excursions \\
        Path variability & Low (paths look similar) & High (jumps vary between runs) \\
        \hline
    \end{tabular}
    \caption{Qualitative comparison of Gaussian and L\'{e}vy noise models.}
    \label{tab:comparison}
\end{table}

In both cases the oscillator continues to respond predominantly near the
natural frequency $\omega_0 = \sqrt{k/m}$ between noise events, demonstrating
that the characteristic response of the system is ignorant of the tail behavior of
the noise.

\section{Discussion}

\subsection{Strengths of the Tool}

The simulation effectively illustrates the mathematical distinction between
Gaussian and L\'{e}vy white noise.
A user can slide $\alpha$ continuously from $2.0$ downward and watch the
trajectory transform from smooth random oscillation to jump-dominated behavior,
making the role of the stability index immediately apparent.
The two-panel layout is particularly useful because the $x(t)$ plot and the
phase portrait highlight complementary aspects of the same trajectory:
the time series shows when jumps occur, while the phase portrait shows the
direction and magnitude of the resulting excursions.
Together they give a more complete picture of the system's behavior than
either plot alone.

The Re-run button adds interactivity that a static figure cannot provide.
Because the solution is stochastic, pressing Re-run at fixed parameters
shows a different realization each time, allowing the user to build intuition
about variability: under Gaussian noise successive runs look nearly identical,
while under L\'{e}vy noise the number, timing, and size of jumps can differ
substantially from run to run.

\subsection{Limitations}

\textbf{Numerical stability at small $\alpha$.} The Euler--Maruyama method
can produce unreliable results when individual noise increments are very large
relative to the current state, which occurs frequently at small $\alpha$ and
large $\sigma$. In these regimes the velocity update adds an increment that
can dwarf the drift term, occasionally causing the trajectory to diverge in
a single step. A smaller time step or an implicit stochastic integrator
would improve robustness in these extreme parameter regimes.

\textbf{Single realization.} The tool only displays one sample path at a time.
Some useful statistics such as the mean displacement
$\mathbb{E}[x(t)]$, the variance $\mathrm{Var}(x(t))$, or the marginal
distribution of $x$ require averaging over many realizations, which a single path cannot show.

\textbf{Scope.} The simulation is limited to a linear oscillator with scalar
additive noise. Extensions would include nonlinear restoring forces
(Van der Pol or Duffing oscillator), multiplicative noise where the noise
amplitude depends on the state, or the display of additional diagnostics
such as the power spectral density, which would reveal the peak at the
natural frequency $f_0$ directly.

\end{document}